\section{\benchname: a blind ground-truth benchmark}
\label{sec:bench}
The core weakness of version 2 was simple: its production claims rested on detectors
that had never been checked against a known answer.  \benchname{} fixes that.  We
serve \emph{known} models behind the same anonymous contract the scanner faces in
production, freeze the protocol, run the scanner, and only then look at the labels.

\subsection{Subjects and axes}
We serve \BenchN{} subjects (Table~\ref{tab:subjects}).  The dense decoders run from
128 to 2{,}048 in residual width: Pythia, OPT, GPT-2, DistilGPT-2, SmolLM, Qwen2.5-0.5B,
and TinyLlama-1.1B.  There are two Mamba SSMs (130M and 370M), two MoE subjects (the
trained OLMoE-1B-7B, and a random-init Mixtral mechanics subject), a random-init
DeepSeek-V3-config MLA mechanics subject, four masked encoders from the BERT family, and
one encoder-decoder (T5-small).  Around this core we vary the \emph{serving}, not the
weights: the same checkpoint computed in fp32, bf16, and fp16; with its wire values
rounded to bf16; with int8 dynamic quantisation; behind an 8-way batching condition;
and a router fixture that alternates two checkpoints behind one model name.  The
random-init subjects exercise real code paths (MoE routing, MLA attention) even though
their weights are meaningless.  We label them mechanics subjects and disable the
behavioural probes for them.

\textbf{Effective sample size.}  The \BenchN{} subjects are \emph{not} \BenchN{}
independent architectures.  They span \BenchArchs{} distinct base checkpoints.  The
serving axis adds five GPT-2 variant configurations (four completed; fp16 omitted) plus a
two-checkpoint router.  Those variants exist to measure serving confounds, not to
inflate $n$.  All headline accuracies therefore carry cluster-bootstrapped intervals at
the base-checkpoint level (Table~\ref{tab:accuracy}, \BenchClusters{} clusters), and we
never quote the variant-inflated count as architectural coverage.  \BenchNScored{} of
\BenchN{} subjects completed the frozen protocol.  The fp16-on-CPU variant ran out of
time budget, because fp16 emulation on a CPU is several times slower than fp32 or bf16.
We report it as omitted rather than quietly dropping it.

\subsection{The anonymous contract}
Every subject answers the same OpenAI-compatible \texttt{/v1/chat/completions} endpoint.
It provides honest usage accounting, top-$k$ logprobs with $k\le20$, \texttt{logit\_bias},
SSE streaming with server-reported per-token times (in the style of Azure's
\texttt{latency\_checkpoint}), a stated context cap, and an explicit
\texttt{measurement\_columns} diagnostic that is used \emph{only} in the rich-access
regime.  Ground truth is recorded from the loaded config classes, never typed by hand.
The scanner sees no headers, no model strings, and no config hints.

\subsection{Blindness and freezing}
Three mechanisms keep the benchmark honest.  (i) The battery protocol, meaning the probe
set, the sample sizes, the spectral thresholds, and the decode-growth decision
thresholds, is frozen in code and hashed (\texttt{protocol \ProtocolHash}) before any
label is analysed.  (ii) The controlled-change study is run commit-reveal.  The harness
draws each episode's label in secret, writes only a SHA-256 commitment to disk before
scanning, and after the scan the reveal is verified against the commitment
(Sec.~\ref{sec:blindfixtures}).  (iii) Detectors that failed during development are
documented in the appendix together with their fixes.  One was a curvature gate that
fired without any effect-size check; another was a units error in the decode-growth
slope.  The benchmark reports the \emph{repaired} instrument, and the failure log is
part of the record.

\subsection{What the benchmark measures}
For each property we report three things on the support scores actually shown in this
paper: top-1 accuracy, the Brier score, and the expected calibration error (ECE).  Top-1
accuracy asks whether the highest-ranked candidate is right.  The Brier score is the
mean squared gap between the score and the truth, so lower is better and 0 is perfect.
ECE asks whether a score of 0.8 is right about 80\% of the time.  All three are reported
under two prior views (the reference prior and the evidence-only flat prior) and one
sensitivity check (the legacy independent-factor combination).  For each numeric target
we report rank-recovery coverage by width and the relative error on resolved estimates,
separately for rich and for production-realistic access.  For each serving axis we
report how the detectors respond to changes in dtype, batching, and routing.  The
baselines (Sec.~\ref{sec:baselines}) run on the same recordings.
